% Capítulo de la tesis (texto derivado de envios/paper_empirico/cuerpo.tex; se edita aquí).
\chapter{Trabajos futuros}
\label{sec:agenda-futura}

La agenda se desprende de lo que el panel no permite establecer.

\begin{enumerate}
\item \textbf{Más economías}, para que la inferencia no dependa de trece \textit{clusters}: completar el EMBI de Indonesia, Filipinas y Sudáfrica entre 2004 y 2023 y el de India antes de 2012, y homologar el $GaR$ de Argentina con el del resto del panel para incorporarla a la muestra principal (como robustez no cambia ninguna conclusión; Sección~\ref{sec:robustez-argentina}).
\item \textbf{Contrastar la heterogeneidad con un criterio ex ante.} La amplificación condicional de la Sección~\ref{sec:robustez-heterogeneidad} usa una división de economías definida a partir de una versión anterior. El paso siguiente es reemplazarla por un moderador continuo y medido de forma independiente ---la participación extranjera en la deuda soberana---, de modo que el contraste deje de depender de la búsqueda de especificación y permita contrastar si la amplificación depende de quién fija el precio de la deuda.
\item \textbf{Variación exógena específicamente bancaria}, como cambios regulatorios o choques de fondeo idiosincrásicos a bancos grandes, que desplacen la fragilidad sin pasar por el spread y permitan identificar un efecto causal.
\item \textbf{Mayor frecuencia}: con datos diarios o semanales de precios bancarios y spreads, la secuencia temporal permitiría separar la dirección de la retroalimentación mejor que con trimestres.
\item \textbf{Una lista ex ante de episodios sin respaldo oficial} más larga que la disponible, para que la prueba sobre el estrés emergente no descanse en un solo episodio de tres trimestres.
\item \textbf{Controles trimestrales y la calificación soberana histórica}, para completar la réplica de \citet{Chari2024b} y separar los canales fiscal y cambiario del efecto de la fragilidad.
\end{enumerate}

Una pregunta adicional es si la concentración bancaria amplifica la complementariedad. La teoría no tiene una predicción unívoca \citep{MMR2010b,BoydDeNicolo2005}, y una triple interacción con tres medidas de concentración, estimada en la especificación en niveles, no identifica ningún efecto.
